This section separates what the descriptive pass suggests from what the inferential stage decides. The first is reported here only to establish that there is something to measure; the second is reported under the rules fixed in \autoref{tab:bayesDecisions}, and returns no verdict it has not earned.

Before the Bayesian estimates are reported, the pages of \autoref{sec:appSweep} provide a descriptive reading of what those models must explain. Across the retrained geometries, dominant-frequency reconstruction appears more reliable over the lower half of the analysed band, while faster inputs often return near $0$\,Hz. Within that lower half, possible holes occur near $\mathcal F_{\mathrm{lock}}$, and hierarchical band probes sometimes weaken there, although the local $f\pm1$\,Hz task remains stronger. The generated-signal examples add a second question: slow components may survive, yet their forecast amplitudes can approach the noise floor, making recovery ambiguous. These plots therefore motivate paired Bayesian analysis on generated backgrounds and frame the contrasts reported next. They are exploratory observations only, not confirmation of structural aliasing or of the preregistered hypotheses in \autoref{sec:four}.
